% year: תשפ"ו
% semester: א
% moed: ב
% number: 1ב
% points: 10
% categories: func-limits
% official_solution: yes
% source: exams/מבחנים/תשפו/sol B.pdf

\begin{question}
חשבו את הגבול הבא
\[ \lim_{x\to\infty} \left( \frac{x^3 + 1}{x^3 - 2} \right)^{x^3 + 3} \]
\end{question}

\begin{hint}
זה ביטוי מהצורה $1^\infty$. כתבו $\frac{x^3 + 1}{x^3 - 2} = 1 + \frac{3}{x^3 - 2}$ והשתמשו בגבול $\left(1 + \frac{1}{u}\right)^u \to e$.
\end{hint}

\begin{solution}
כאשר $x \to \infty$ הבסיס שואף ל-$1$ והמעריך ל-$\infty$ (ביטוי $1^\infty$). נכתוב
\[ \frac{x^3 + 1}{x^3 - 2} = \frac{(x^3 - 2) + 3}{x^3 - 2} = 1 + \frac{3}{x^3 - 2}, \qquad x^3 + 3 = (x^3 - 2) + 5, \]
ולכן (עבור $x$ גדול, כך ש-$x^3 - 2 > 0$)
\[ \left( \frac{x^3 + 1}{x^3 - 2} \right)^{x^3 + 3} = \left( 1 + \frac{3}{x^3 - 2} \right)^{x^3 - 2} \cdot \left( 1 + \frac{3}{x^3 - 2} \right)^{5}. \]
נסמן $u = \frac{x^3 - 2}{3} \to \infty$. אז הגורם הראשון שווה $\left[\left(1 + \frac{1}{u}\right)^{u}\right]^3 \to e^3$, לפי הגבול היסודי $\lim_{u\to\infty}\left(1 + \frac1u\right)^u = e$ ורציפות $t \mapsto t^3$.
הגורם השני שואף ל-$1^5 = 1$, כי $\frac{3}{x^3 - 2} \to 0$.
לפי אריתמטיקה של גבולות:
\[ \boxed{\lim_{x\to\infty} \left( \frac{x^3 + 1}{x^3 - 2} \right)^{x^3 + 3} = e^3 \cdot 1 = e^3}. \]
(דרך חלופית: $\ln$ של הביטוי הוא $(x^3 + 3)\ln\left(1 + \frac{3}{x^3 - 2}\right)$, ומכיוון ש-$\ln(1 + t) \sim t$ כאשר $t \to 0$, הוא שואף ל-$\lim \frac{3(x^3 + 3)}{x^3 - 2} = 3$.)
\end{solution}
